% Generated draft: requires mathematical and source review.
\documentclass[11pt]{article}
\usepackage[margin=25mm]{geometry}
\usepackage{fontspec}
\setmainfont{Latin Modern Roman}
\usepackage{amsmath,amssymb,mathtools,unicode-math}
\setmathfont{Latin Modern Math}
\usepackage{enumitem,needspace,xurl,hyperref}
\setlength{\parindent}{0pt}
\setlength{\parskip}{5pt}
\setlength{\emergencystretch}{3em}
\newcommand{\N}{\mathbb{N}}
\newcommand{\Z}{\mathbb{Z}}
\newcommand{\Q}{\mathbb{Q}}
\newcommand{\R}{\mathbb{R}}
\newcommand{\C}{\mathbb{C}}
\newcommand{\F}{\mathbb{F}}
\newcommand{\A}{\mathbb{A}}
\newcommand{\PP}{\mathbb{P}}
\newcommand{\E}{\mathbb{E}}
\newcommand{\eps}{\varepsilon}
\newcommand{\dd}{\,\mathrm{d}}
\begin{document}
\section*{UnsolvedMath ID: 30002598}
\textbf{Source problem code:} catalogue entry
\subsection{Definitions and mathematical statement}
A presimplicial module consists of modules \(C_n\), indexed by integers \(n\geq 0\), over a coefficient ring \(k\), together with \(k\)-linear face maps \(d_i^{(n)}:C_n\to C_{n-1}\) for \(n\geq 1\) and \(0\leq i\leq n\), satisfying
\[
 d_i^{(n-1)}\circ d_j^{(n)}=d_{j-1}^{(n-1)}\circ d_i^{(n)}
 \qquad(n\geq 2,\ 0\leq i<j\leq n).
\]
Its alternating boundary in positive degrees is
\[
 \partial_n=\sum_{i=0}^{n}(-1)^i d_i^{(n)}.
\]
The source does not specify the coefficient ring, module-sidedness, or an augmentation in degree zero.

Let \(q\) denote a formal indeterminate. The source proposes the expressions
\[
 C_n^q=C_n\otimes_k\mathbb Z[q],\qquad
 \partial_n^q=\sum_{i=0}^{n}q^i\bigl(d_i^{(n)}\otimes 1\bigr):C_n^q\longrightarrow C_{n-1}^q
 \quad(n\geq 1).
\]
These expressions are conditional on compatible coefficient structures making the tensor product, induced face maps, and multiplication by \(q\) well-defined. Such structures are not specified in the record: in particular, \(\mathbb Z[q]\) is not automatically a \(k\)-algebra for an arbitrary coefficient ring \(k\). The displayed tensor product is retained exactly as supplied, without replacing it by a different scalar extension.

\textbf{Exploratory problem.} For a given presimplicial module, and subject to resolving the coefficient issue above, investigate whether the weighted operators \(\partial_n^q\) yield a useful homological structure. The source explicitly allows the general failure of
\[
 \partial_{n-1}^q\circ\partial_n^q=0\qquad(n\geq 2).
\]
Thus ordinary chain-complex homology is not automatically defined by these operators. The original wording asks whether this construction is useful; the normalized wording asks what useful homological structure can be extracted. Neither specifies a proposed invariant, axioms, a class of modules, restrictions or specializations on \(q\), or a criterion of usefulness. Consequently, the supplied problem is an exploratory research question rather than a fully specified proposition with a fixed conclusion.
\subsection{Short English statement}
Given a presimplicial module, investigate whether weighting its face maps by \(1,q,q^2,\ldots\) produces a useful homological structure, even though the resulting operators generally do not square to zero. The coefficient extension \(C_n\otimes_k\mathbb Z[q]\) requires unspecified compatibility data, and the intended meaning of a useful homological structure remains open to clarification.
\subsection{Sources}
\begin{itemize}
\item Supplied local UnsolvedMath registry record, permanent ID \(30002598\), title \emph{Weighted \(q\)-Boundaries for Presimplicial Modules}, collection \emph{Oberwolfach Reports Open Problems}; exact source problem code \texttt{catalogue entry}. The record gives proposed year 2014 and identifies no individual proposer.
\item Supplied report citation: \emph{Algebraic Structures in Low-Dimensional Topology} (2015), Oberwolfach Reports, pp.\ 1403--1458, DOI \texttt{10.4171/owr/2014/26}. \url{https://doi.org/10.4171/owr/2014/26}. The supplied citation's year 2015 differs from the report designation 26/2014 and proposed year 2014; this discrepancy is unresolved here.
\item Supplied literature source: \emph{Oberwolfach Report 26/2014}. \url{https://publications.mfo.de/bitstream/handle/mfo/3416/OWR_2014_26.pdf?isAllowed=y&sequence=1}.
\item Supplied literature source: Przytycki, \emph{Progress in distributive homology}. \url{https://arxiv.org/abs/1406.6499}. This reference does not identify Przytycki as the proposer of this particular problem.
\item Supplied literature assessment, checked \texttt{2026-08-21}: the registry labels the problem open, describes the report's question as exploratory, notes that the weighted map is not generically a chain differential, and reports that no authoritative source was found settling a useful general theory for the construction. This is supplied dated curation derived from the OWR material, not an independently verified literature assessment; the status and cited claims require independent verification.
\end{itemize}
\end{document}
